\documentclass[11pt]{article}
\usepackage[T1]{fontenc}
\usepackage{lmodern,amsmath,amssymb,amsthm,booktabs,array,longtable}
\usepackage[margin=27mm]{geometry}
\usepackage[colorlinks=true,linkcolor=blue,citecolor=blue,urlcolor=blue]{hyperref}
\newtheorem{theorem}{Theorem}
\newtheorem{proposition}[theorem]{Proposition}
\newtheorem{corollary}[theorem]{Corollary}
\newcommand{\NN}{\mathbb N}
\newcommand{\ZZ}{\mathbb Z}
\title{Resolving the double-root conjecture in\protect\\
Ballot--Kimberling--Moses (2017),\protect\\
Section 5, Table 5, p.~23}
\author{Haobo Ma\\\small ChronoAI Pte Ltd, Singapore
\and Wenlin Zhang\\\small National University of Singapore\\
\small ChronoAI Pte Ltd, Singapore}
\date{3 November 2026}
\hypersetup{pdftitle={Resolving the double-root conjecture: Ballot--Kimberling--Moses (2017), Section 5, Table 5, p. 23},
pdfauthor={Haobo Ma and Wenlin Zhang}}
\begin{document}
\maketitle
\begin{abstract}
We prove the conjectured recurrence for $q(x)=x^2-2x+1$ in
Ballot--Kimberling--Moses, Section 5, Table 5, p.~23 (2017).
The polynomial-form tree has generation sizes $G_n=1+\binom n2$ for
every $n\geq0$, so its recurrence with coefficients $(3,-3,1)$ holds
from $n=3$. Coordinates given by value and derivative at $1$ give an
exact shortest-path formula. The same counting approach also proves
the $x^2-1$ row's eventual recurrence and corrects its printed value
$G_3=4$ to $G_3=3$. A hexagonal-lattice argument proves the neighboring
$x^2-x+1$ row: $G_n=6n-20$ for $n\geq6$, with its recurrence starting
at $n=8$. We also count the generations modulo $x^m-1$ for every $m\geq1$.
\end{abstract}

\section{The polynomial-form tree}
Ballot, Kimberling and Moses define the polynomial-form tree $T(q)$ by
reducing the operations $p\mapsto p+1$ and $p\mapsto xp$ modulo a fixed
monic polynomial $q$, starting from $0$ and removing duplicates
\cite[Section 1, pp.~15--17]{bkm}. Their Section 5 describes the table
recurrences as conjectured eventual recurrences. We treat the
double-root row of Table 5, where $q=(x-1)^2$.
The rows $x^2-1$ and $x^2-x+1$ admit comparably short distance proofs;
we give both below. For context, the $x^2-2$ case was proved in
Stoll's 2015 answer \cite{stoll}, and the $x^2-3$ case is covered by
\cite[Theorem 4.3]{bkm}. For the golden-ratio row $x^2-x-1$,
OEIS A252864 records Wiktorowicz's proof in August 2026
\cite{golden,wiktorowicz}. These are prior proofs of different rows.
The double-root recurrence also has the pinned Lean proof described
in the final section.

For precision, work in $R=\ZZ[x]/((x-1)^2)$. Let $S_n$ be the set of
residues reachable by at most $n$ operations. Thus $S_0=\{0\}$ and
\[
 S_{n+1}=S_n\cup\{p+1:p\in S_n\}\cup\{xp:p\in S_n\}.
\]
Set $g_0=\{0\}$, $g_n=S_n\setminus S_{n-1}$ for $n\geq1$, and
$G_n=|g_n|$. This agrees with first-occurrence deletion in the source:
replacing a duplicate prefix by an earlier occurrence cannot delay
any subsequent residue, since both operations respect congruence.
Choosing a different parent among duplicates at the same depth changes
edges, not generation sets.

Every residue has a unique representative
\[
 p=a+b(x-1),\qquad (a,b)=(p(1),p'(1))\in\ZZ^2.
\]
Indeed, writing $y=x-1$ and reducing modulo $y^2$ gives exactly its
constant and linear coefficients. In these coordinates the operations are
\begin{equation}\label{eq:operations}
 A(a,b)=(a+1,b),\qquad M(a,b)=(a,b+a).
\end{equation}
The reachable pairs are $(0,0)$ and all $(a,b)$ with $a\geq1$, $b\geq0$,
as the proof below shows. These are polynomial residues, not just their
values at $1$: evaluation forgets $b$ and gives the different numeric
tree $T(1)$.

\section{Exact distances and the recurrence}
\begin{theorem}\label{thm:main}
For $a\geq1$ and $b\geq0$, the first generation of $a+b(x-1)$ is
\begin{equation}\label{eq:distance}
 d(a,b)=a+\left\lceil\frac ba\right\rceil.
\end{equation}
The root has distance $0$. For $n\geq1$, the complete generation is
\begin{equation}\label{eq:shell}
 g_n=\{(n,0)\}\ \dot\cup\!
 \bigcup_{a=1}^{n-1}\{(a,b):a(n-a-1)<b\leq a(n-a)\}.
\end{equation}
Consequently, for every $n\geq0$,
\begin{equation}\label{eq:count}
 G_n=1+\binom n2,
\end{equation}
and for every $n\geq3$,
\begin{equation}\label{eq:recurrence}
 G_n=3G_{n-1}-3G_{n-2}+G_{n-3}.
\end{equation}
Here $\binom02=\binom12=0$.
\end{theorem}
\begin{proof}
An operation word reaching $(a,b)$ contains exactly $a$ additions.
If it has $m$ multiplications, each increases the second coordinate by
an integer between $0$ and $a$. Hence $b\leq am$, and its length is
at least $a+\lceil b/a\rceil$. Nonnegativity follows from
\eqref{eq:operations}; if $a=0$, every operation is a multiplication of
the root and $b=0$.

For the reverse bound, divide $b=at+s$ with $t\geq0$ and $0\leq s<a$.
If $s=0$, apply $A$ $a$ times, then $M$ $t$ times. If $s>0$, apply,
in order, $A$ $s$ times, $M$ once, $A$ $a-s$ times, and $M$ $t$ times.
The endpoints are respectively $(a,at)$ and $(a,s+at)$, and the lengths
are respectively $a+t$ and $a+t+1$. This proves \eqref{eq:distance}
and reachability for every stated pair.

For $d(a,b)=n$, either $a=n$ and $b=0$, or $1\leq a<n$ and
$\lceil b/a\rceil=n-a$. The latter equality is exactly the interval
in \eqref{eq:shell}; it contains $a$ integers. The intervals have
different first coordinates, so
$G_n=1+\sum_{a=1}^{n-1}a=1+n(n-1)/2$ for $n\geq1$.
The root gives the same formula at $n=0$. Its third finite difference
vanishes, proving \eqref{eq:recurrence}.
\end{proof}

This gives $G_0,\ldots,G_{10}=1,1,2,4,7,11,16,22,29,37,46$, exactly
the double-root row. The counts from $n=1$ are the classical sequence
A000124 with index shifted by one \cite{oeis}; the tree identification
is the content of the distance argument.

The accumulated sets have the equally explicit description
\[
 S_n=\{(0,0)\}\ \dot\cup\!
 \bigcup_{a=1}^n\{(a,b):0\leq b\leq a(n-a)\},\qquad
 |S_n|=1+n+\binom{n+1}{3}.
\]
The set formula follows directly from \eqref{eq:distance}; the cardinality
follows by summing $a(n-a)+1$ over $1\leq a\leq n$.
Also, as a formal power series,
\[
 \sum_{n\geq0}G_nz^n
 =\frac1{1-z}+\frac{z^2}{(1-z)^3}
 =\frac{1-2z+2z^2}{(1-z)^3}.
\]
The series identity follows by expanding $(1-z)^{-3}$, whose coefficient
of $z^{n-2}$ is $\binom n2$.

\section{The neighboring row \texorpdfstring{$q(x)=x^2-1$}{q(x)=x2-1}}
\begin{proposition}\label{prop:sibling}
In $\ZZ[x]/(x^2-1)$, use the same two operations and first-occurrence
generations. Then $G_0=1$ and $G_n=n$ for every $n\geq1$.
In particular, $G_n=2G_{n-1}-G_{n-2}$ for all $n\geq3$.
The Table 5 initial entry $G_3=4$ for this row is $3$ instead.
\end{proposition}
\begin{proof}
Represent a residue uniquely by $ux+v$. The operations on $(u,v)$ are
$(u,v)\mapsto(u,v+1)$ and $(u,v)\mapsto(v,u)$.
Starting at $(0,0)$, all coordinates are nonnegative; each addition
increases $u+v$ by one and a multiplication preserves it.
If $u=0$, the distance to $(0,v)$ is $v$. If $u>0$, at least one
multiplication is necessary, so the distance is at least $u+v+1$.
Adding $u$ times, multiplying once, and adding $v$ times attains this
bound. Thus generation $n\geq1$ consists of $(0,n)$ together with
the $n-1$ pairs satisfying $u\geq1$, $v\geq0$, $u+v=n-1$.
The recurrence follows from $G_n=n$ at indices $n,n-1,n-2\geq1$.
At depth $3$ the residues are exactly $3,2x,x+1$: the candidate $x^2$
equals the already present residue $1$.
\end{proof}

\begin{corollary}\label{cor:cyclic}
For the polynomial-form tree modulo $x^m-1$, where $m\geq1$, write a
reachable residue uniquely as $\sum_{j=0}^{m-1}c_jx^j$ with $c_j\geq0$.
Every such residue is reachable. For a nonzero residue, let
$r=\max\{j:c_j>0\}$ and $C=\sum c_j$. Its distance is $C+r$.
Consequently $G_0=1$ and
\[
 G_n=\sum_{r=0}^{m-1}\binom{n-1}{r}\quad(n\geq1),\qquad
 \sum_{n\geq0}G_nz^n=1+\sum_{r=0}^{m-1}\frac{z^{r+1}}{(1-z)^{r+1}}.
\]
Here a binomial coefficient is zero when $r>n-1$.
In particular, $m=3$ has the same generation sizes as the double-root
tree, at every depth, although the quotient polynomials differ.
\end{corollary}
\begin{proof}
Multiplication cyclically rotates the $m$ coefficients and addition
increases their sum by one. Thus any word has exactly $C$ additions.
If it has fewer than $r$ multiplications, every added unit has final
exponent less than $r$, even after reduction modulo $m$; this cannot
produce $c_r>0$. Horner's construction with $r$ multiplications and
$C$ additions attains $C+r$. For fixed $r$, the condition $C=n-r$
with $c_r\geq1$ has $\binom{C+r-1}{r}=\binom{n-1}{r}$ solutions,
by distributing $C-1$ units into $r+1$ positions. Summing over $r$
gives the count and then the displayed formal generating function.
\end{proof}

\section{The hexagonal row \texorpdfstring{$q(x)=x^2-x+1$}{q(x)=x2-x+1}}
\begin{proposition}\label{prop:hex}
For $q=x^2-x+1$, the first eleven generation sizes are
\[
G_0,\ldots,G_{10}=1,1,2,4,7,11,16,22,28,34,40.
\]
More precisely, $G_0=1$ and, for $n\geq1$,
\begin{equation}\label{eq:hexcount}
G_n=\sum_{j=0}^{5}\mathbf1_{n\geq j+1}
       +\sum_{j=1}^{5}(n-j-1)_+ +(n-6)_+,
\end{equation}
where $t_+=\max(t,0)$. Hence $G_n=6n-20$ for $n\geq6$ and
$G_n=2G_{n-1}-G_{n-2}$ for every $n\geq8$; this onset is sharp.
\end{proposition}
\begin{proof}
Identify $a+bx$ with $(a,b)$. Since $x^2=x-1$ and $x^6=1$, the
successive powers of $x$ are the six rays generated by
\[
 \begin{aligned}
 e_0&=(1,0),& e_1&=(0,1),& e_2&=(-1,1),\\
 e_3&=(-1,0),& e_4&=(0,-1),& e_5&=(1,-1).
 \end{aligned}
\]
Every nonzero integer pair lies either on exactly one ray $t e_j$
($t\geq1$) or strictly between two successive rays, in the form
$s e_j+t e_{j+1}$ ($s,t\geq1$), with indices taken modulo $6$.
Indeed the six cones partition the plane and every adjacent determinant
is $1$, so the coefficients are integers and unique. On each cone,
\[
 L(a,b)=\max\{|a|,|b|,|a+b|\}
\]
is the sum of those coefficients. This is a norm and $L(e_j)=1$.

A word with $h$ additions and $k$ multiplications represents its
endpoint as a sum of $h$ vectors from $e_0,\ldots,e_k$, with subscripts
reduced modulo $6$. Therefore $h\geq L(a,b)$.
Conversely a sum using powers through $x^k$ can be formed by Horner's
construction in $h+k$ steps.
For rays $e_j$ with $0\leq j\leq4$, a nonzero multiple cannot lie in
the cone generated by the preceding rays. The same holds for the
strict sector between $e_j$ and $e_{j+1}$, for $0\leq j\leq3$,
when only rays through $e_j$ are allowed. Thus the exact distances are
$t+j$ on these rays and $s+t+j+1$ in these four strict sectors:
the cone restriction and norm give the lower bounds, and the adjacent
representations give attaining words.

The remaining ray $t e_5$ and the two remaining strict sectors have
negative second coordinate. At least four multiplications are needed,
because $e_0,\ldots,e_3$ have nonnegative second coordinate. If there
are at least five, the norm bound gives length at least $L+5$.
With exactly four, write the endpoint either as $t e_5$, or as
$s e_4+t e_5$, or as $s e_5+t e_0$. In each case the linear functional
$F(a,b)=a-b$ has value at least $L+1$, whereas
$F(e_j)\leq1$ for $0\leq j\leq4$. Hence at least $L+1$ additions
are needed, again giving length at least $L+5$.
The adjacent representations with five multiplications attain this
bound. We have therefore obtained $t+5$ on the last ray and
$s+t+5$ in each of the last two strict sectors.

At depth $n$, ray $j$ contributes $\mathbf1_{n\geq j+1}$.
A strict sector of offset $r$ contributes $(n-r-1)_+$, since
$s+t=n-r$ with $s,t\geq1$. The six offsets are $1,2,3,4,5,5$,
proving \eqref{eq:hexcount}. Its terms sum to $6n-20$ for $n\geq6$.
The recurrence follows for $n\geq8$; at $n=7$ its discrepancy is
$22-2\cdot16+11=1$, proving sharpness.
\end{proof}

\begin{corollary}
For the hexagonal row the formal generating function is
\[
 \sum_{n\geq0}G_nz^n=
 \frac{1-z+z^2+z^3+z^4+z^5+z^6+z^7}{(1-z)^2}.
\]
\end{corollary}
\begin{proof}
Multiply the series by $(1-z)^2$. The coefficient at $n$ is
$G_n-2G_{n-1}+G_{n-2}$, taking negative subscripts as zero.
Proposition~\ref{prop:hex} gives the displayed coefficients through
$n=7$ and zero thereafter.
\end{proof}

\section{Coverage and verification}
The main theorem settles the individual double-root row, not a family
obtained by varying its polynomial. The source's numbered Conjectures
1 and 6 have parameters strictly beyond the boundary $k=2$; the present
result does not remove those restrictions or prove those families.
The following table accounts individually for the rows of Tables 5 and 6
and the numbered conjectures. ``Not covered'' records the scope of
these arguments, rather than a claim that the item is still open.
The $x^2+1$ example in Section 3 already has the source's written
count $G_n=4n-9$ for $n\geq4$. The cross-referenced MathOverflow
question \cite{stoll} has Stoll's 2015 proof for $\sqrt2$ and
Wiktorowicz's August 2026 answer for the golden ratio. Its broader
question about other $r$ is not settled by these examples.
The stronger strict inequality suggested in Meyerowitz's answer
\cite{meyerowitz} has a small counterexample: for the $\sqrt2$ tree,
$\ell(2\sqrt2)=\ell(3)=3$. Both endpoints have a three-step word,
and neither is in $\{0,1,2,\sqrt2\}$, the complete set reached in at
most two steps. This is the equality case $j=1$, $b=0$ of that answer's
comparison $\ell(b+(j+1)\sqrt2)\leq\ell(2j+1+b\sqrt2)$.
\small
\begin{longtable}{@{}p{0.42\linewidth}p{0.52\linewidth}@{}}
\toprule
Source item in \cite{bkm} & Coverage here\\
\midrule
\endhead
Table 5, $x^2-2x+1$, p.~23 & Proved: all generation sets and counts,
recurrence from $n=3$.\\
Table 5, $x^2-1$, p.~23 & Refuted-and-corrected: printed $G_3=4$
is $3$; proved $G_n=n$ for $n\geq1$.\\
Table 5, $x^2-x+1$ & Proved: Proposition~\ref{prop:hex}, recurrence from $8$.\\
Table 5, $x^2-2$ & Prior proof cited: Stoll (2015); source Theorem 4.3.\\
Table 5, $x^2-3$ & Prior proof cited: source Theorem 4.3.\\
Table 5, $x^2-x-1$ & Prior proof: Wiktorowicz \cite{wiktorowicz},
$G_n=G_{n-1}+G_{n-3}$ for $n\geq12$.\\
Table 5, $x^2-4$ & Not covered.\\
Table 5, $x^2-3x+1$ & Not covered.\\
Table 5, $x^2-4x+1$ & Not covered.\\
Table 5, $x^2-5x+1$ & Not covered.\\
Table 5, $x^2-3x-1$ & Not covered.\\
Table 5, $x^2-5x-1$ & Not covered.\\
Table 5, $x^2-2x-1$ & Not covered.\\
Table 5, $x^2-4x-1$ & Not covered.\\
Table 5, $x^2-6x-1$ & Not covered.\\
Table 6, $x^2-3x+2$, p.~24 & Not covered.\\
Table 6, $x^2-4x+2$ & Not covered.\\
Table 6, $x^2-5x+2$ & Not covered.\\
Table 6, $x^2-6x+2$ & Not covered.\\
Table 6, $x^2-4x+3$ & Not covered.\\
Table 6, $x^2-5x+3$ & Not covered.\\
Table 6, $x^2-6x+3$ & Not covered.\\
Table 6, $x^2-7x+3$ & Not covered.\\
Conjecture 1, $x^2-kx+1$, $k>2$ & Not covered.\\
Conjecture 2, $x^2-(2k+1)x-1$, $k>0$ & Not covered.\\
Conjecture 3, $x^2-2kx-1$, $k>1$ & Not covered.\\
Conjecture 4, $x^2-kx+2$, $k>3$ & Not covered.\\
Conjecture 5, $x^2-kx+3$, $k>4$ & Not covered.\\
Conjecture 6, $x^2-kx+(k-1)$, $k>2$ & Not covered.\\
Section 1, Theorem 1.1 & Prior proof cited: the source's universal tree.\\
Section 3, $T(i)$ & Prior proof cited: $G_n=4n-9$ for $n\geq4$.\\
Section 4, Theorem 4.3 and composition remarks & Prior proofs cited in the source;
no new composition result here.\\
Section 6, Theorems 6.1--6.2 and Table 7 & Not covered; the source states
these recurrences.\\
MathOverflow 195207, Stoll's general-$r$ expectation & Not covered.\\
MathOverflow answer 195264, strict predecessor inequality & Refuted-and-corrected:
$\ell(2\sqrt2)=\ell(3)=3$.\\
\bottomrule
\end{longtable}
\normalsize

The \href{https://github.com/the-omega-institute/trureturing/blob/50f2a97bb5e704842f2c32c6095f68a0dc62c779/D5/S3/Combinatorics/PolynomialTrees/BallotKimberlingMosesDoubleRootGeneration.lean}{pinned Lean source}
defines the children by monic polynomial remainders and generations by
successive accumulated-set differences. In the namespace
\begin{center}\small
\texttt{D5.S3.Combinatorics.PolynomialTrees.}\\
\texttt{BallotKimberlingMosesDoubleRootGeneration}
\end{center}
the declaration \texttt{result : claim} proves the eventual recurrence
\eqref{eq:recurrence}, with witness $N=3$.
The source-local lemma \texttt{seen\_spec} establishes the reachable sets;
\texttt{generation\_succ\_spec} establishes the generation intervals;
\texttt{G\_sum} establishes the triangular-sum count.
The shortest-path formula, generating functions, accumulated-set
cardinality, Propositions~\ref{prop:sibling} and \ref{prop:hex}, and
Corollary~\ref{cor:cyclic} have written proofs here.
The accompanying \texttt{verify.py} checks exact polynomial remainders,
pair coordinates, complete generation sets, distances, counts, and
recurrences through depth $80$ for the three quadratic polynomials.
It checks the cyclic-quotient formulas for $1\leq m\leq6$ through depth $16$.
These finite checks
are separate from the proofs for arbitrary depth.

\begin{thebibliography}{9}
\bibitem{bkm}
C.~Ballot, C.~Kimberling and P.~J.~C.~Moses,
\emph{Linear Recurrences Originating From Polynomial Trees},
The Fibonacci Quarterly \textbf{55}(5) (2017), 15--27.
\href{https://doi.org/10.1080/00150517.2017.12427731}{doi:10.1080/00150517.2017.12427731}.
\href{https://www.fq.math.ca/Papers1/55-5/BallotKimberlingMoses.pdf}{Journal text}.
\bibitem{oeis}
OEIS Foundation Inc., \emph{The On-Line Encyclopedia of Integer Sequences},
\href{https://oeis.org/A000124}{A000124}, accessed 3 November 2026.
\bibitem{stoll}
M.~Stoll, answer to C.~Kimberling,
\emph{A possibly surprising appearance of Lucas numbers},
MathOverflow, 30 January 2015.
\href{https://mathoverflow.net/questions/195207}{Question and answers}.
\bibitem{meyerowitz}
A.~Meyerowitz, answer to C.~Kimberling,
\emph{A possibly surprising appearance of Lucas numbers},
MathOverflow, 30 January 2015.
\href{https://mathoverflow.net/a/195264}{Answer 195264}.
\bibitem{golden}
OEIS Foundation Inc.,
\href{https://oeis.org/A252864}{A252864}, proof sketch by
R.~Wiktorowicz first posted 14 August 2026
(\href{https://oeis.org/history/view?seq=A252864\&v=35}{revision 35}),
with revised proof text dated 24 August 2026;
accessed 3 November 2026.
\bibitem{wiktorowicz}
R.~Wiktorowicz, \emph{oeis-lean: machine-checked proofs of OEIS
conjectures in Lean 4}, version 1.0.0, 25 August 2026.
\href{https://doi.org/10.5281/zenodo.22092442}{doi:10.5281/zenodo.22092442}.
\end{thebibliography}
\noindent\small\copyright\ 2026 Haobo Ma and Wenlin Zhang.
Licensed under \href{https://creativecommons.org/licenses/by/4.0/}{Creative Commons Attribution 4.0 International (CC BY 4.0)}.
\end{document}
